\chapter{A circle becomes a sheet}\label{connection:sheets}
\regionline{Thickness, surface and the hand}
\atlasfiguretitle{A circle becomes a sheet: paired comparison}
\comparisonplate{Pita}{Arepa}{A pocket made by puffing beside a thick round opened with a knife.}

Put a tortilla, chapati and kisra on the table and the first likeness is easy to see. All spread a small amount of mixture across a broad cooking surface. Their thinness keeps the heat close to the middle. Once they are cooked, however, the hand finds three different structures. The tortilla begins as maize dough, the chapati as wheat dough and kisra as fermented sorghum batter. A circle is a meeting place for different materials.

For the tortilla recipe, nixtamalised masa harina takes up water to become a dough that can be pressed. Chapati's atta is kneaded and rested before rolling. Kisra must flow under the spreading tool. The cook who tries to improve all three by adding flour until nothing feels sticky would lose what each mixture is meant to do. Their recipes begin with an ingredient and a desired texture, rather than with a single ideal of dough.\citep{patitortilla08,manaliroti06,moggakisra}

\section{The work before the pan}
Salgado's tortilla demonstration starts farther back, with maize kernels cooked and steeped in an alkaline solution before grinding. Masa harina makes that preparation available as a purchased ingredient in the home recipe. The packet is convenient because earlier work has already happened. It is not an invitation to replace it with ordinary cornmeal.\citep{tortilla,cimmytnixt08}

The arepa recipe begins with another prepared maize product, masarepa. It absorbs water, rests and becomes a thicker round. Broa begins with maize flour scalded for a loaf containing other flours. Set the three packets beside one another and maize stops being a single interchangeable shelf item. Treatment, grind and the bread being made matter as much as the crop's name.\citep{goyaarepa08,broarecipe}

This is one of the pleasures of making breads from different regions. The shopping becomes part of the learning. Read the ingredient label, make the first small portion and watch what water does to it. A bread can teach the cook why an apparently minor distinction on a bag deserves attention.

\section{A pocket, a layer, a fold}
Pita's puff can be spectacular. A well-rolled disc meets a hot surface, its walls separate and a pocket appears. Yet not every puffed bread is meant to keep that space for filling. Pane carasau also swells, then the baker splits it into sheets and returns them to the oven to dry. One empty space becomes a pocket, the other becomes access to two thin pieces. The next action makes the difference.\citep{aboodpita,carasaurecipe}

A thick arepa can be cut to make room for a filling, so its pocket comes from the knife rather than inflation. Nothing about the sandwich requires these two breads to share a dough history. Their comparison is more useful when it stays with the task: how does the cook create space, and how does the finished wall hold a meal?

Lavash and yufka open another possibility. A sufficiently broad, flexible sheet can fold around food without a pocket at all. A stack is convenient because its layers can be separated as needed. UNESCO's flatbread-sharing account describes preparations involving different ovens and griddles, as well as collective making and sharing. The cooking surface belongs in the comparison because it changes how a sheet is handled from dough to finished bread.\citep{lavash,sharing}

\section{Thin does not always mean plain}
Msemen and lachha paratha turn thinness inward. Their cooks stretch or roll dough, add fat and fold or coil the sheet before cooking. What looks like a compact portion contains several layers. Tear the warm bread and the work inside becomes visible. The recipes ask for different shaping sequences, but both reward patient handling before the bread reaches the pan.\citep{benmsemen,rakslachha06}

Shaobing's paste supplies a further comparison. Here a layer carries flavour as well as separation. Recco focaccia carries cheese between sheets, while the printed shaobing uses a seasoned sesame mixture. The filling changes what escapes when the bread is opened and what the cook needs to keep inside during shaping.\citep{linliushaobing07,reccofocaccia}

Appam brings us back to batter. Its curved pan lets a thin rim surround a deeper centre. Thickness varies within one bread, so the crisp edge and softer middle can share a plate. The pan contributes a form that a flat cooking surface would not produce in the same way. It is worth looking at the equipment before assuming that an unfamiliar recipe is merely another version of a familiar disc.\citep{familiarappam06}

A flatbread is therefore a beginning for comparison, not the end of one. Its edge, centre, folds and cooking marks tell us what happened before it reached the table. Looking at those small features makes distant likenesses more rewarding, because each bread keeps its own method in view.

\section{What belongs inside the dough}
A disc's ingredients can also explain a likeness that its outline leaves hidden. In Sri Lankan pol roti, coconut is distributed through the dough. Break the round and the shreds appear in the middle as well as at its toasted surface. The bread carries some of its richness before an accompaniment is added.\citep{roshanipol06}

Dini's account follows that preparation into a different kitchen after leaving Sri Lanka. Fresh coconut became harder to obtain, so she used unsweetened desiccated coconut with coconut milk. The adaptation begins with an understanding of what the ingredient contributes. It preserves the wish to make the bread without pretending that every shop offers the same materials.\citep{dinipol06}

The Sri Lankan food guide places wheat-flour pol roti beside kurakkan roti made with finger millet, brown rice flour and coconut. The round's name and use can remain familiar while the grain changes its flavour and handling. The comparison therefore belongs both within Sri Lanka and beside the chapati elsewhere in the atlas.\citep{srilankafoods06}

Tear each bread before adding sauce or a filling. The face gives the first impression, but the broken edge often explains the preparation more clearly. A layer, a coconut shred, a pocket wall or a dense maize crumb can show why the cook mixed, rested or shaped the portion as they did. A table of flat breads becomes more interesting when the interior is allowed to speak too.

\begin{breadnote}{Make space in two ways}
Bake pita on page~68 and arepas on page~146. Open one of each before filling. Compare the walls, the crumb and the amount of filling they comfortably hold. For a sheet that changes purpose after puffing, turn to carasau on page~27.
\end{breadnote}
